\documentclass[11pt]{article}
\usepackage[a4paper,margin=26mm]{geometry}
\usepackage{amsmath,amssymb}
\usepackage[hidelinks]{hyperref}
\setlength{\emergencystretch}{3em}
\title{Conjecture 00000003890: infinitely many jumps of one fixed ideal}
\author{lizaixi01, with OpenAI Codex assistance}
\date{22 October 2026}
\begin{document}
\maketitle
\section{A fixed ideal and the definition}
The conjecture claims that the full set of F-jumping parameters of a
fixed ideal is always finite. Let $R=\mathbf F_2[X]$ and fix the nonzero
proper ideal $I=(X)$. Every positive integer is an F-jumping parameter
of $I$. Thus the first clause is false; no assertion about the second
clause, concerning unions inside $(0,1)$, is needed. We count all positive
jumps, as the source states, rather than restricting to a bounded interval.

For $q=2^s$ with $s\geq1$, the Frobenius power $J^{[q]}$ is the ideal
generated by the $q$th powers of elements of $J$. The Frobenius root
$A^{[1/q]}$ is the smallest ideal $J$ with $A\subseteq J^{[q]}$.
For $t\geq0$, the generalized test ideal is
\[
 \tau(I^t)=\bigcup_{s\geq1}(I^{\lceil tq\rceil})^{[1/q]},
 \qquad q=2^s.                                          \tag{1}
\]
This is the standard definition in the regular F-finite polynomial-ring
setting; see Definition 4.1 of \cite{BC}. The ring $\mathbf F_2[X]$ is
within that setting. Below we compute actual ideals, and verify that the
family in (1) is increasing, so its union agrees with its ideal supremum.

\section{The least Frobenius root is computed}
For every $N\geq0$,
\[
 ((X^N))^{[1/q]}=(X^{\lfloor N/q\rfloor}).                \tag{2}
\]
Indeed, the right-hand ideal has Frobenius power
$(X^{q\lfloor N/q\rfloor})$, which contains $(X^N)$.
For minimality, let $J=(g)$ be any ideal whose Frobenius power contains
$X^N$. The polynomial ring is a principal ideal domain, and
$g^q\mid X^N$. Since $X$ is prime, $g$ is associated to $X^j$ for some
$j\geq0$. The divisibility gives $qj\leq N$, hence
$j\leq\lfloor N/q\rfloor$. Therefore
$(X^{\lfloor N/q\rfloor})\subseteq J$, proving minimality against every
candidate ideal, not only against other monomial ideals.

Formula (2) identifies the $s$th ideal in (1) as
$(X^{\lfloor\lceil tq\rceil/q\rfloor})$.
The inequality $\lceil2tq\rceil\leq2\lceil tq\rceil$ shows that these
ideals increase with $s$. They decrease when $t$ increases.
At a nonnegative integer $k$, every ideal in (1) is $(X^k)$, giving
\[
 \tau(I^k)=(X^k).                                       \tag{3}
\]

\section{Every positive integer is a jump}
For $k\geq0$ and $q=2^s$, define
\[
 t_{k,s}=\frac{kq+q-1}{q}=k+1-\frac1q.
\]
Then $0\leq t_{k,s}<k+1$ and $t_{k,s}\to k+1$ from the left.
The ideal in (1) at index $s$ is $(X^k)$, since
$\lceil t_{k,s}q\rceil=kq+q-1$ and
$\lfloor(kq+q-1)/q\rfloor=k$. Consequently
\[
 X^k\in\tau(I^{t_{k,s}}),\qquad
 X^k\notin\tau(I^{k+1})=(X^{k+1}).
\]
For any $0\leq t<k+1$, choose $s$ large enough that $t<t_{k,s}$.
The decreasing dependence of $\tau(I^t)$ on $t$ then implies
$\tau(I^t)\ne\tau(I^{k+1})$. This proves the usual left-jump criterion
at every $k+1$, and in particular rules out local constancy from the left.
The injective map $k\mapsto k+1$ from $\mathbf N$ into the jump set proves
that this fixed ideal has infinitely many jumps.

\section{Formalization and verification}
The Lean ring is the actual \texttt{Polynomial (ZMod 2)}. Frobenius powers
use \texttt{Ideal.map} of \texttt{iterateFrobenius}. The root is an
\texttt{sInf} over all candidate ideals, and
\texttt{frobeniusRoot\_\allowbreak X\_\allowbreak pow\_\allowbreak least}
proves that the infimum is attained and is the least candidate.
Thus existence or minimality of roots is not assumed.

\texttt{testStages\_\allowbreak monotone} proves the increasing chain,
and \texttt{mem\_\allowbreak testIdeal\_\allowbreak iff} identifies the
ideal supremum with the literal union. \texttt{testIdeal\_\allowbreak integer}
and \texttt{testIdeal\_\allowbreak antitone} prove (3) and parameter
monotonicity. The left-approaching points and their monomial witnesses
are explicit. \texttt{positive\_\allowbreak integer\_\allowbreak jump}
uses the criterion $\tau(I^t)\ne\tau(I^c)$ for every $0\leq t<c$.
Finally \texttt{infinitely\_\allowbreak many\_\allowbreak jumps} is a
\texttt{Set.Infinite} theorem about the real jump parameters.

The pinned project uses Lean 4.33.0 and exact Mathlib/dependency commits.
Independent project build, fifteen-declaration axiom audit and official
fresh-kernel replay are recorded. No admitted proof, native decision
shortcut or extra axiom is used. The original bilingual source is fixed at
upstream commit\par\noindent
\texttt{b355e42906e6eaf52bc7e9d005e0ef02d8aa7400}.\par
OpenAI Codex assisted the proof and report; a separate supervising Codex
agent reviewed the original, definitions and full formal chain. This is
AI review, not human certification. The example is classical.

\begin{thebibliography}{1}
\bibitem{BC} W. Badilla-C\'espedes and E. Le\'on-Cardenal,
\emph{F-Pure Thresholds and F-Volumes of Some Non-Principal Ideals},
Mediterranean Journal of Mathematics 22, 98 (2025), Definition 4.1.
\url{https://arxiv.org/html/2305.00571v3}.
\end{thebibliography}
\end{document}
